% !TEX program = xelatex
% 论文精读入口。修改 config.tex 后，在本文件填写实际论文内容。
\documentclass[aspectratio=169,10pt,fontset=none]{ctexbeamer}
\usetheme{BITUniversal}
% 所有入口共用。复制项目后，先改这里的标题、姓名、机构和日期。
\title[组会汇报]{组会汇报}
\subtitle{研究主题 / 论文中文标题}
\setenglishtitle{Research Topic / Original Paper Title}
\setpresenter{汇报人姓名}
\setaffiliation{北京理工大学}
\setmeeting{课题组组会}
\date{\number\year.\ifnum\month<10 0\fi\number\month.\ifnum\day<10 0\fi\number\day}
% \date{2026.11.09}  % 如需固定日期，取消本行注释。
\setpaper{Original Paper Title}{Author A, Author B, et al.}{Conference / Journal, Year}

% 默认：DT 参考中的轻量白底页眉；改成 band 使用 PWBC 绿色渐变页眉。
\setheaderstyle{light}
% \bitsectionpagesfalse  % 短汇报可关闭自动章节过渡页。

\begin{document}
\BITtitlepage
\BITcontents
\bitsection{背景与相关工作}{研究问题、已有方法与本文定位}
\begin{frame}{论文信息与核心问题}
  \begin{block}{\BITpaper}
    \BITauthors\\[.5em]{\small\color{BITgray}\BITvenue}
  \end{block}
  \begin{itemize}
    \item \bitplaceholder{研究对象、应用条件与现有瓶颈。}
    \item \bitplaceholder{相对最相关方法的关键改进。}
  \end{itemize}
  \bitcallout{核心贡献：\bitplaceholder{一句话概括最有价值的技术改进。}}
\end{frame}
\bitsection{方法与训练流程}{整体框架、关键设计与运行方式}
\begin{frame}{整体框架}
  \bitimageplaceholder[.56\textheight]{插入论文框架图}
  \bitsource{来源：论文 Fig. X。}
\end{frame}
\begin{frame}{关键设计与训练推理}
  \bitcolumns{\begin{block}{关键机制}
    \bitplaceholder{输入、输出与设计动机。}\end{block}}%
    {\begin{block}{训练与推理}
    \bitplaceholder{训练目标、推理条件与计算开销。}\end{block}}
  \bitcallout{关键作用：\bitplaceholder{说明设计如何解决前面的问题。}}
\end{frame}
\bitsection{实验结果与分析}{主要结果、消融与适用范围}
\begin{frame}{实验设置与主要结果}
  \bitexperiment{\bitplaceholder{验证核心贡献。}}{\bitplaceholder{数据、基线、指标与统计口径。}}
  \bitimageplaceholder[.43\textheight]{插入结果图或替换为原生表格}
  \bitcallout{结果：\bitplaceholder{关键数值、比较对象及结论。}}
\end{frame}
\begin{frame}{消融实验与现象解释}
  \bitexperiment{\bitplaceholder{验证关键模块的作用。}}{\bitplaceholder{保持其余条件一致。}}
  \bitimageplaceholder[.43\textheight]{插入消融结果或定性对比}
  \bitcallout{解释：\bitplaceholder{结合设计动机解释结果。}}
\end{frame}
\bitsection{总结与讨论}{技术结论与本课题的联系}
\begin{frame}{总结与课题启发}
  \begin{itemize}
    \item \bitplaceholder{最重要的贡献与证据。}
    \item \bitplaceholder{已验证的范围及局限。}
    \item \bitplaceholder{可复用的设计和需要额外验证的假设。}
  \end{itemize}
\end{frame}
\BITthankyou
\end{document}
